% ============================================================================
% LANGENTWURF LE-04d: Wiederholung + Stundenleistung - Mi colegio
% Spanisch Klasse 7 - Sequenz 4: Mi colegio
% ============================================================================
% Tier: 1 (vollständig)
% Doppelstunde: 4d (Std. 7-8 von 8) - ABSCHLUSS
% Fachfarbe: #c41e3a (Karminrot)
% ============================================================================

\documentclass[11pt,a4paper]{article}

\usepackage[utf8]{inputenc}
\usepackage[T1]{fontenc}
\usepackage[ngerman]{babel}
\usepackage{geometry}
\usepackage{fancyhdr}
\usepackage{lastpage}
\usepackage{xcolor}
\usepackage{tcolorbox}
\usepackage{enumitem}
\usepackage{tabularx}
\usepackage{longtable}
\usepackage{array}
\usepackage{booktabs}
\usepackage{amssymb}

% ============================================================================
% KONFIGURATION
% ============================================================================

\newcommand{\fach}{Spanisch}
\newcommand{\fachfarbe}{c41e3a}
\newcommand{\jahrgangsstufe}{7}
\newcommand{\schulart}{Gesamtschule}
\newcommand{\stundenthema}{Wiederholung + Stundenleistung: Mi colegio}
\newcommand{\reihenthema}{Mi colegio}
\newcommand{\stundennummer}{4d}
\newcommand{\reihenstunden}{4}
\newcommand{\gession}{A1}
\newcommand{\lehrwerk}{Qué pasa 1}
\newcommand{\lehrwerkseite}{Unidad 2}
\newcommand{\lerngruppengroesse}{24}

% ============================================================================
% LAYOUT
% ============================================================================
\geometry{left=2.5cm, right=2.5cm, top=2.5cm, bottom=2.5cm}
\definecolor{fach}{HTML}{\fachfarbe}
\definecolor{gray}{HTML}{666666}
\definecolor{lightgray}{HTML}{f5f5f5}
\definecolor{successgreen}{HTML}{2a7a4b}
\definecolor{warningorange}{HTML}{b35c00}
\definecolor{basis}{HTML}{2a7a4b}
\definecolor{standard}{HTML}{2c5aa0}
\definecolor{erweiterung}{HTML}{7b1fa2}

\tcbuselibrary{skins,breakable}

\newtcolorbox{infobox}[1][]{
    colback=lightgray,
    colframe=fach,
    fonttitle=\bfseries,
    title=#1,
    breakable
}

\newtcolorbox{scorebox}{
    colback=white,
    colframe=fach,
    fonttitle=\bfseries,
    title=Didaktik-Score,
    breakable
}

\pagestyle{fancy}
\fancyhf{}
\fancyhead[L]{\small\textcolor{gray}{\fach{} -- Jg. \jahrgangsstufe{} -- \stundenthema}}
\fancyhead[R]{\small\textcolor{gray}{LE-\stundennummer{} | Tier 1}}
\fancyfoot[C]{\small\thepage{} / \pageref{LastPage}}
\renewcommand{\headrulewidth}{0.4pt}

% Differenzierungsmarker
\newcommand{\B}{\textcolor{basis}{\textbf{[B]}}}
\newcommand{\Std}{\textcolor{standard}{\textbf{[S]}}}
\newcommand{\E}{\textcolor{erweiterung}{\textbf{[E]}}}

% ============================================================================
% DOKUMENT
% ============================================================================
\begin{document}

% ============================================================================
% DECKBLATT
% ============================================================================
\begin{titlepage}
    \centering
    \vspace*{2cm}
    {\large\textcolor{gray}{\schulart}}\\[2cm]
    {\Huge\textcolor{fach}{\textbf{Langentwurf}}}\\[0.5cm]
    {\large LE-\stundennummer{} | Tier 1 (vollständig)}\\[2cm]
    {\LARGE\textbf{\stundenthema}}\\[0.5cm]
    {\large Sequenz: \reihenthema{} -- \textbf{ABSCHLUSS}}\\[2cm]
    \begin{tabular}{rl}
        \textbf{Fach:} & \fach{} (A1) \\[0.3cm]
        \textbf{Klasse:} & \jahrgangsstufe \\[0.3cm]
        \textbf{Lehrwerk:} & \lehrwerk, \lehrwerkseite \\[0.3cm]
        \textbf{Doppelstunde:} & \stundennummer{} (Std. 7-8 von 8) \\[0.3cm]
    \end{tabular}
    \vfill
\end{titlepage}

\tableofcontents
\newpage

% ============================================================================
% 1. BEDINGUNGSANALYSE
% ============================================================================
\section{Bedingungsanalyse}

\subsection{Lerngruppe}
\begin{tabular}{ll}
    Kursgröße: & \lerngruppengroesse{} SuS \\
    Jahrgangsstufe: & \jahrgangsstufe \\
    Sprachniveau: & A1 (GER) -- Anfänger \\
    Fremdsprache: & 2. Fremdsprache (nach Englisch) \\
\end{tabular}

\subsection{Vorwissen aus 04a/04b/04c}
\begin{itemize}
    \item \textbf{04a:} Schulfächer (las asignaturas): matemáticas, inglés, español, música, historia, educación física, arte, ciencias naturales
    \item \textbf{04b:} Wochentage (los días de la semana) + Uhrzeiten (¿A qué hora...? A las...)
    \item \textbf{04c:} Verben auf -ar konjugieren: estudiar, hablar, cantar, escuchar, terminar
\end{itemize}

\subsection{Differenzierung (intern)}
\begin{tabular}{|l|p{3.5cm}|p{3.5cm}|p{3.5cm}|}
\hline
& \B{} \textbf{Basis} & \Std{} \textbf{Standard} & \E{} \textbf{Erweiterung} \\
\hline
\textbf{Ziel} & BR: Rezeption, Grundwortschatz & MR: Produktion mit Hilfe & GY: Freie Anwendung \\
\hline
\textbf{Wortschatz} & 6-8 Schulfächer & 10+ Schulfächer & + Meinungsäußerung \\
\hline
\textbf{-ar Verben} & Lückentext einsetzen & Sätze bilden & Freie Texte schreiben \\
\hline
\textbf{Test} & Basisaufgaben (60\%) & Alle Pflichtaufgaben & + Zusatzaufgaben \\
\hline
\end{tabular}

% ============================================================================
% 2. SACHANALYSE
% ============================================================================
\section{Sachanalyse}

\subsection{Sprachliche Strukturen}

\textbf{Schulfächer (las asignaturas):}
\begin{center}
\begin{tabular}{|l|l|}
\hline
\textbf{Spanisch} & \textbf{Deutsch} \\
\hline
las matemáticas & Mathematik \\
el inglés & Englisch \\
el español & Spanisch \\
la música & Musik \\
la historia & Geschichte \\
la educación física (E.F.) & Sport \\
el arte & Kunst \\
las ciencias naturales & Naturwissenschaften \\
la informática & Informatik \\
la geografía & Erdkunde \\
\hline
\end{tabular}
\end{center}

\textbf{Wochentage (los días de la semana):}
\begin{center}
\begin{tabular}{|l|l|l|l|l|l|l|}
\hline
lunes & martes & miércoles & jueves & viernes & sábado & domingo \\
\hline
Montag & Dienstag & Mittwoch & Donnerstag & Freitag & Samstag & Sonntag \\
\hline
\end{tabular}
\end{center}

\textbf{Uhrzeiten:}
\begin{itemize}
    \item ¿A qué hora...? = Um wie viel Uhr...?
    \item A la una = Um ein Uhr
    \item A las dos / tres / cuatro... = Um zwei / drei / vier Uhr...
    \item y media = halb, y cuarto = viertel nach
\end{itemize}

\textbf{Verben auf -ar (Konjugation):}
\begin{center}
\begin{tabular}{|l|l|l|l|l|}
\hline
\textbf{Person} & \textbf{estudiar} & \textbf{hablar} & \textbf{cantar} & \textbf{escuchar} \\
\hline
yo & estudio & hablo & canto & escucho \\
tú & estudias & hablas & cantas & escuchas \\
él/ella & estudia & habla & canta & escucha \\
nosotros & estudiamos & hablamos & cantamos & escuchamos \\
(vosotros) & (estudiáis) & (habláis) & (cantáis) & (escucháis) \\
ellos/ellas & estudian & hablan & cantan & escuchan \\
\hline
\end{tabular}
\end{center}

\textbf{Endungen:} -o, -as, -a, -amos, (-áis), -an

% ============================================================================
% 3. DIDAKTISCHE ANALYSE
% ============================================================================
\section{Didaktische Analyse}

\subsection{Stellung in der Sequenz}
Dies ist die \textbf{4. und letzte Doppelstunde} (Std. 7-8 von 8) der Sequenz ``\reihenthema''.
Die Stunde schließt die Sequenz ab mit:
\begin{enumerate}
    \item Konsolidierung aller bisherigen Inhalte (Schulfächer, Wochentage, Uhrzeiten, -ar Verben)
    \item Wiederholungsübungen zur Vorbereitung
    \item Stundenleistung (summative Bewertung)
\end{enumerate}

\subsection{Kompetenzziele (Rahmenplan MV)}
\begin{itemize}
    \item \textbf{Sprechen:} Über den Schulalltag berichten
    \item \textbf{Schreiben:} Kurze Texte über den eigenen Stundenplan
    \item \textbf{Verfügen über sprachliche Mittel:} -ar Verben konjugieren, Wortschatz Schule
    \item \textbf{Interkulturell:} Schulsystem in Spanien vs. Deutschland
\end{itemize}

\subsection{Legitimation}
\textbf{Rahmenplan Spanisch MV:}
\begin{itemize}
    \item An Gesprächen teilnehmen: über den Schulalltag sprechen
    \item Verfügen über sprachliche Mittel: Regelmäßige Verben konjugieren
\end{itemize}

% ============================================================================
% 4. STUNDENZIELE
% ============================================================================
\section{Stundenziele}

\subsection{Grobziel}
Die SuS \textbf{konsolidieren} die Inhalte der Sequenz 4 (Schulfächer, Wochentage, Uhrzeiten, -ar Verben) und weisen ihre Kompetenzen in einer \textbf{Stundenleistung} nach.

\subsection{Feinziele (operationalisiert)}
\begin{tabular}{|c|p{8cm}|c|c|}
\hline
\textbf{Nr.} & \textbf{Die SuS können...} & \textbf{AFB} & \textbf{Diff.} \\
\hline
FZ1 & mindestens 8 Schulfächer auf Spanisch benennen und zuordnen & I & alle \\
\hline
FZ2 & die Wochentage korrekt anwenden und Uhrzeiten angeben & I/II & alle \\
\hline
FZ3 & regelmäßige -ar Verben in allen Personen konjugieren & II & alle \\
\hline
FZ4 & einen kurzen Text über ihren Schulalltag verfassen & II/III & \Std{}/\E{} \\
\hline
\end{tabular}

% ============================================================================
% 5. METHODISCHE ANALYSE
% ============================================================================
\section{Methodische Analyse}

\subsection{Medien}
\begin{itemize}
    \item Tafel/Whiteboard
    \item Lehrwerk: \lehrwerk, \lehrwerkseite
    \item Arbeitsblatt: AB-04d.pdf (Wiederholung)
    \item Stundenleistung: SL-04d.pdf (Test)
    \item Optional: LP-04d.html (Wiederholung)
\end{itemize}

\subsection{Sozialformen}
\begin{itemize}
    \item Plenum: Schnelle Wiederholung, Fragen klären
    \item Partnerarbeit: Dialogübung (Stundenplan)
    \item Einzelarbeit: Wiederholungsübungen + Stundenleistung
\end{itemize}

\subsection{Besonderheit: Stundenleistung}
\begin{itemize}
    \item Dauer: 25 Minuten
    \item Umfang: 20 Punkte (4 Aufgaben)
    \item Bewertung: Note nach Prozentskala
    \item Differenzierung: Basis 60\% = ausreichend
\end{itemize}

% ============================================================================
% 6. VERLAUFSPLANUNG
% ============================================================================
\section{Verlaufsplanung}

\textbf{1. Stunde (45 Min.) -- Schwerpunkt: Wiederholung}

\begin{longtable}{|p{1cm}|p{1.5cm}|p{4cm}|p{3cm}|p{2cm}|p{2cm}|}
\hline
\textbf{Zeit} & \textbf{Phase} & \textbf{Lehrerhandeln} & \textbf{Schülerhandeln} & \textbf{SF} & \textbf{Medien} \\
\hline
\endhead

0--5 & Einstieg & ``¿Qué asignatura tenéis ahora?'' Schnelle Aktivierung & Antworten & Plenum & -- \\
\hline

5--15 & Wiederhol. & Schulfächer-Memory an der Tafel & Zuordnen, Nachsprechen & Plenum & Tafel \\
\hline

15--25 & Übung 1 & AB-04d Aufg. 1-2 (Schulfächer, Wochentage) & Bearbeiten & EA & AB \\
\hline

25--35 & Übung 2 & -ar Verben Konjugationswettbewerb (Teams) & Konjugieren & Teams & Tafel \\
\hline

35--45 & Dialog & PA: ``¿Qué tienes los lunes?'' & Dialogübung & PA & AB \\
\hline
\end{longtable}

\vspace{1em}

\textbf{2. Stunde (45 Min.) -- Schwerpunkt: Stundenleistung}

\begin{longtable}{|p{1cm}|p{1.5cm}|p{4cm}|p{3cm}|p{2cm}|p{2cm}|}
\hline
\textbf{Zeit} & \textbf{Phase} & \textbf{Lehrerhandeln} & \textbf{Schülerhandeln} & \textbf{SF} & \textbf{Medien} \\
\hline
\endhead

0--5 & Aktivierung & Letzte Fragen klären, Mut machen & Fragen stellen & Plenum & -- \\
\hline

5--10 & Vorbereitung & SL erklären, Zeitansage & Zuhören & Plenum & SL \\
\hline

10--35 & \textbf{TEST} & Aufsicht, individuelle Hilfe bei Verständnisfragen & Stundenleistung bearbeiten & EA & SL \\
\hline

35--40 & Abgabe & Einsammeln & Abgeben & -- & -- \\
\hline

40--45 & Abschluss & Ausblick Sequenz 5, Feedback & Rückmeldung & Plenum & -- \\
\hline
\end{longtable}

\textbf{Legende:} EA = Einzelarbeit, PA = Partnerarbeit, SF = Sozialform, SL = Stundenleistung

% ============================================================================
% 7. ERWARTETE SCHWIERIGKEITEN
% ============================================================================
\section{Erwartete Schwierigkeiten}

\begin{tabular}{|p{5cm}|p{8cm}|}
\hline
\textbf{Schwierigkeit} & \textbf{Reaktion} \\
\hline
Verwechslung -ar Endungen & Konjugationstabelle an der Tafel \\
\hline
Artikel bei Schulfächern & Merkregel: la música, el inglés, las matemáticas (Plural!) \\
\hline
Uhrzeiten (a la/a las) & A la una (1), a las dos/tres... (Rest) \\
\hline
Prüfungsangst & Ruhige Atmosphäre, klare Zeitansagen, \B{} Zeitzugabe \\
\hline
\end{tabular}

% ============================================================================
% 8. HAUSAUFGABE
% ============================================================================
\section{Hausaufgabe}

\begin{itemize}
    \item \textbf{Vor der Stunde:} LP-04d.html zur Wiederholung (empfohlen)
    \item \textbf{Nach der Stunde:} Vokabeln Sequenz 4 wiederholen für Vokabeltest
\end{itemize}

% ============================================================================
% 9. LÖSUNGEN
% ============================================================================
\section{Lösungen AB-04d}

\textbf{Aufgabe 1 (Schulfächer zuordnen):} (4P)
\begin{itemize}
    \item las matemáticas = Mathematik
    \item la música = Musik
    \item la historia = Geschichte
    \item la educación física = Sport
\end{itemize}

\textbf{Aufgabe 2 (Wochentage + Uhrzeiten):} (4P)
\begin{itemize}
    \item lunes = Montag, martes = Dienstag, miércoles = Mittwoch, jueves = Donnerstag
    \item A las ocho = Um 8 Uhr, A las nueve y media = Um halb 10
\end{itemize}

\textbf{Aufgabe 3 (-ar Verben konjugieren):} (6P)
\begin{itemize}
    \item estudiar: estudio, estudias, estudia, estudiamos, estudian
    \item hablar: hablo, hablas, habla, hablamos, hablan
\end{itemize}

\textbf{Aufgabe 4 (Schulalltag beschreiben):} (6P)
\begin{itemize}
    \item Beispiel: Los lunes tengo matemáticas a las ocho. Estudio español los martes. Me gusta la música.
\end{itemize}

% ============================================================================
% 10. STUNDENLEISTUNG - ÜBERSICHT
% ============================================================================
\section{Stundenleistung SL-04d}

\textbf{Aufbau:} 4 Aufgaben, 20 Punkte, 25 Minuten

\begin{tabular}{|c|l|c|c|c|}
\hline
\textbf{Aufg.} & \textbf{Inhalt} & \textbf{AFB} & \textbf{Punkte} & \textbf{Diff.} \\
\hline
1 & Zuordnung: Schulfächer & I & 4 & alle \\
\hline
2 & Wochentage + Uhrzeiten & I/II & 4 & alle \\
\hline
3 & -ar Verben konjugieren & II & 6 & alle \\
\hline
4 & Schulalltag beschreiben & II/III & 6 & alle \\
\hline
\multicolumn{3}{|r|}{\textbf{Gesamt:}} & \textbf{20} & \\
\hline
\end{tabular}

\textbf{Bewertung (Klasse 7):}
\begin{tabular}{|c|c|c|c|c|c|c|}
\hline
Note & 1 & 2 & 3 & 4 & 5 & 6 \\
\hline
ab \% & 96 & 80 & 60 & 40 & 20 & <20 \\
\hline
ab Punkte & 19 & 16 & 12 & 8 & 4 & <4 \\
\hline
\end{tabular}

% ============================================================================
% 11. DIDAKTIK-SCORE
% ============================================================================
\newpage
\section{Didaktik-Score}

\begin{scorebox}
\textbf{Bewertung der didaktischen Qualität nach 10 Dimensionen}\\
\textbf{Ziel-Score: $\geq$ 9.0 (Premium-Qualität)}

\vspace{1em}
\begin{tabular}{|c|l|c|c|p{4.5cm}|}
\hline
\textbf{\#} & \textbf{Dimension} & \textbf{Gew.} & \textbf{Score} & \textbf{Notiz} \\
\hline
1 & Differenzierung & 15\% & 9/10 & SL mit Basis-Sicherung, Zeitzugabe \\
\hline
2 & Taxonomie (AFB) & 10\% & 9/10 & AFB I: 40\%, AFB II: 45\%, AFB III: 15\% \\
\hline
3 & Lernziel-Alignment & 10\% & 10/10 & Rahmenplan: Schulalltag \\
\hline
4 & Aktivierung & 10\% & 9/10 & Memory, Wettbewerb, Dialoge \\
\hline
5 & Scaffolding & 10\% & 9/10 & Wiederholung $\rightarrow$ Übung $\rightarrow$ Test \\
\hline
6 & Feedback-Qualität & 10\% & 9/10 & Bewertungsraster transparent \\
\hline
7 & Sprachliche Klarheit & 10\% & 10/10 & Altersgerecht Kl. 7 \\
\hline
8 & Motivation \& Relevanz & 10\% & 9/10 & Alltagsrelevanz: eigener Stundenplan \\
\hline
9 & Inklusion (UDL) & 10\% & 9/10 & Zeitzugabe, mündlich/schriftlich \\
\hline
10 & Praktikabilität & 5\% & 9/10 & 45+45 Min realistisch \\
\hline
\multicolumn{3}{|r|}{\textbf{GESAMT:}} & \textbf{9.1/10} & \textcolor{successgreen}{\textbf{FREIGABE}} \\
\hline
\end{tabular}
\end{scorebox}

\end{document}
